\documentclass[../../../include/open-logic-section]{subfiles}
%READERNOTE{SOL6-A246: सर्व संक्रमक संचांत विभक्तीकरण सत्य असते हा मूळ दावा चुकीचा आहे. आवश्यक उपसंच-समावेशकत्वाची अट सिद्धतेच्या लघुतम प्रतिमानात आणि सरावाच्या विधानात जोडली. सरावाची सिद्धता दिलेली नाही. सांत स्वयंसिद्धकीकरणीयतेच्या पहिल्या निष्कर्षात सुसंगततेची अट आणि प्रतिमानांच्या अरिक्ततेचा आधार स्पष्ट केला.}

\begin{document}

\olfileid{sth}{replacement}{finiteaxiomatizability}
\olsection[सांत स्वयंसिद्धकीकरणीयता]{परिशिष्ट: सांत स्वयंसिद्धकीकरणीयता}
प्रतिस्थापनापासून काही निष्कर्ष काढून हे प्रकरण संपवू. पहिला निष्कर्ष \citet{Montague1961} यांचा आहे. लक्षात घ्या: हा $\ZF$ च्या \emph{अंतर्गत} पुरावा नाही, तर $\ZF$ \emph{विषयीचा} पुरावा आहे:
\begin{thm}\ollabel{zfnotfinitely}
  $\ZF$ सुसंगत असेल, तर ती सांत स्वयंसिद्धकांनी मांडता येत नाही. अधिक सर्वसाधारणपणे: $\Th{T}$ सांत असेल आणि $\Th{T} \Proves \ZF$ असेल, तर $\Th{T}$ विसंगत असते.

  (येथे आपण गृहीत धरतो की विचाराधीन प्रथम-क्रम वाक्यांमध्ये $\in$ हेच एकमेव अतार्किक मूलचिन्ह आहे; तसेच प्रत्येक $\phi \in \ZF$ साठी $\Th{T} \Proves \phi$ आहे हे दाखवण्यासाठी $\Th{T} \Proves \ZF$ असे लिहितो.)
\end{thm}
\begin{proof}
  $\Th{T} \Proves \ZF$ पूर्ण करणारी सांत $\Th{T}$ निश्चित करा. म्हणून $\Th{T}$ प्रतिबिंबन, म्हणजे \olref[sth][replacement][ref]{reflectionschema}, सिद्ध करते. $\Th{T}$ सांत असल्याने तिचे एका संधीकरणाच्या रूपात, $\theta$ असे, पुनर्लेखन करता येते. $\Th{T}$ अनंतता सिद्ध करत असल्याने, अनंततेचे स्वयंसिद्धकही $\theta$ मध्ये संधीघटक म्हणून घेऊ; त्यामुळे $\theta^X$ पूर्ण करणारा $X$ रिक्त नसतो. या सूत्रावर प्रतिबिंबन वापरल्यास $\Th{T} \Proves \exists \beta(\theta \liff \theta^{V_\beta})$. स्पष्टपणे $\Th{T} \Proves \theta$ असल्याने $\Th{T} \Proves \exists \beta\ \theta^{V_\beta}$ मिळते.

  आता $\psi(X)$ हे पुढीलचे संक्षिप्त लेखन असू द्या:
  \[
    \theta^X \land X\text{ संक्रमक आहे} \land (\forall Y \in X)(Y\text{ संक्रमक आहे}\lif \lnot \theta^{Y})
  \]
साधारणपणे याचा अर्थ: $X$ हे $\theta$ चे संक्रमक प्रतिमान आहे आणि या बाबतीत $\in$-लघुतम आहे. आता $\Th{T} \Proves \exists \beta\ \theta^{V_\beta}$ हे आठवा. $\ZF$ मधील, आणि त्यामुळे $\Th{T}$ मधील, प्रतांकांविषयीच्या मूलभूत तथ्यांवरून:
  \begin{equation}
    \Th{T} \Proves \exists M \psi(M). \tag{*}\label{Mpsi}
  \end{equation}
$\psi(X)$ चे पहिले संधीघटक वापरल्यास, $\Th{T} \Proves \sigma$ असेल तेव्हा $\Th{T} \Proves \forall X(\psi(X) \lif \sigma^X)$ मिळते. म्हणून \eqref{Mpsi} वरून:
  \begin{align*}
    \Th{T} &\Proves \forall X(\psi(X) \lif (\exists N \psi(N))^X)\\
  \intertext{हे आणि पुन्हा \eqref{Mpsi} वापरल्यास:}
    \Th{T} &\Proves \exists M(\psi(M) \land (\exists N \psi(N))^M)
  \intertext{म्हणून विशेषतः:}
    \Th{T} &\Proves \exists M(\psi(M) \land (\exists N \in M)((N\text{ संक्रमक आहे})^M \land (\theta^N)^M))
  \intertext{आता संक्रमणाविषयीच्या प्राथमिक युक्तिवादावरून:}
    \Th{T} &\Proves \exists M(\psi(M) \land (\exists N \in M)(N\text{ संक्रमक आहे} \land \theta^N))
  \end{align*}
म्हणून $\Th{T}$ विसंगत आहे.\footnote{या ``प्राथमिक युक्तिवादा''त संक्रमक संचांसाठी काही ``निरपेक्षतेविषयीची तथ्ये'' सिद्ध करावी लागतात.}
\end{proof}

\citet[223]{Potter2004} यांनी नोंदवलेला पुढील निष्कर्ष यासारखाच आहे:

\begin{prop}\ollabel{finiteextensionofZ}
  $\Th{T}$ ही $\Z$ मध्ये सांत नवीन स्वयंसिद्धके जोडून मिळालेली उपपत्ती असू द्या. $\Th{T} \Proves \ZF$ असेल, तर $\Th{T}$ विसंगत आहे. (येथे \olref{zfnotfinitely} प्रमाणेच गृहीत निर्बंध आहेत.)
\end{prop}
\begin{proof}
  विभक्तीकरणाची (अनंत अनेक) उदाहरणे \emph{वगळून} $\Th{T}$ च्या सर्व स्वयंसिद्धकांचे संधीकरण $\theta$ ने दर्शवा. या पुराव्यासाठी $\psi_*(X)$ हे पुढीलचे संक्षिप्त लेखन असू द्या:
\[
  \theta^X \land X\text{ संक्रमक व उपसंच-समावेशक आहे}
  \land (\forall Y \in X)(Y\text{ संक्रमक व उपसंच-समावेशक आहे}\lif\lnot\theta^Y).
\]
उपसंच-समावेशकत्वाची व्याख्या \olref[sth][spine][Valphabasic]{sec} मध्ये दिली आहे. प्रतिबिंबनातून मिळणारे टप्पे या दोन्ही अटी पूर्ण करतात. म्हणून \olref{zfnotfinitely} च्या प्रतांक-युक्तिवादात या अटी वापरून $\Th{T}\Proves\exists M\psi_*(M)$ मिळते.

  \olref{zfnotfinitely} प्रमाणे, $\Th{T} \Proves \sigma$ असेल तेव्हा $\Th{T} \Proves \forall X(\psi_*(X) \lif \sigma^X)$ मिळते हा रूपबंध सिद्ध करता येतो. मग \olref{zfnotfinitely} प्रमाणेच पुरावा पूर्ण होतो; येथे संक्रमक व उपसंच-समावेशक $M$ मध्ये आतील $N$ साठी या दोन्ही अटींची निरपेक्षता वापरतो.

  तथापि, हा रूपबंध सिद्ध करण्यासाठी \olref{zfnotfinitely} पेक्षा थोडे अधिक काम लागते. कारण विभक्तीकरणाची उदाहरणे $\Th{T}$ मध्ये असली, तरी ती $\theta$ ची संधीघटके नाहीत. पण प्रत्येक विभक्तीकरण-उदाहरण $\sigma$ साठी $\Th{T} \Proves \forall X(X\text{ संक्रमक व उपसंच-समावेशक आहे} \lif \sigma^X)$ हे सिद्ध करून हा अडथळा दूर करता येतो. हे वाचकांसाठी सोडतो.
\end{proof}
\begin{prob}
  प्रत्येक विभक्तीकरण-उदाहरण $\sigma$ साठी $\Z \Proves \forall X(X\text{ संक्रमक व उपसंच-समावेशक आहे} \lif \sigma^X)$ आहे हे दाखवा. (हा रूपबंध आपण \olref[sth][replacement][finiteaxiomatizability]{finiteextensionofZ} मध्ये वापरला.)
\end{prob}
\begin{prob}
  प्रत्येक $\phi \in \Z$ साठी $\ZF \Proves \phi^{V_{\omega+\omega}}$ आहे हे दाखवा.
\end{prob}
\begin{prob}
  \olref[sth][replacement][finiteaxiomatizability]{zfnotfinitely} आणि \olref[sth][replacement][finiteaxiomatizability]{finiteextensionofZ} यांच्या पुराव्यांत वापरलेले उर्वरित रूपबंधात्मक निष्कर्ष तपासा.
\end{prob}

\olref[sth][replacement][absinf]{sec} मध्ये म्हटल्याप्रमाणे, यावरून प्रतिस्थापन हे \olref[ordinals][ordtype]{thmOrdinalRepresentation} पेक्षा काटेकोरपणे अधिक सामर्थ्यवान आहे असे दिसते. किंवा थोडे अधिक नेमके सांगायचे तर: $\Z$ + ``प्रत्येक सुक्रमण एका अद्वितीय क्रमसूचक संख्येशी समरूप आहे'' ही उपपत्ती सुसंगत असेल, तर प्रतिस्थापनाचे काही उदाहरण तिला सिद्ध करता येत नाही.

  %  गृहीतकानुसार $\psi^{V_\alpha}$ चे रूप असे आहे:
  %  \[
  %    (\forall A \in V_\alpha)(\exists S \in V_\alpha)(\forall x \in V_\alpha)(x \in S \liff (\phi^{V_\alpha}(x) \land x \in A))
  %  \]
  %  हे सिद्ध करण्यासाठी $A \in V_\alpha$ निश्चित करा. विभक्तीकरणावरून मिळते:
  %  \[
  %    S = \Setabs{x \in A}{\phi^{V_\alpha}(x)}
  %  \]
  %  आता $S \subseteq A \in V_\alpha$ असल्याने $S \in V_\alpha$; आणि स्पष्टपणे $(\forall x \in V_\alpha)(x \in S \liff (\phi^{V_\alpha}(x) \land x \in A)$.


\end{document}
